%=====================================================================
\chapter{R and Tidyverse Quick Reference}\label{app:r}
%=====================================================================
\normalfigures

R remains the better tool for exploratory statistics, and the department's existing
Quarto lab materials use it. This appendix gives the equivalent of Appendix~A.

\section{Setup}

\begin{lstlisting}[language=R]
install.packages(c("tidyverse", "tidymodels", "GGally", "janitor", "here"))
library(tidyverse)     # dplyr, ggplot2, tidyr, readr, purrr
library(tidymodels)    # recipes, parsnip, rsample, yardstick
set.seed(42)           # reproducibility -- Chapter 29
\end{lstlisting}

\section{Loading and Inspecting}

\begin{lstlisting}[language=R]
df <- read_csv("data.csv")
df <- janitor::clean_names(df)      # snake_case column names

dim(df); glimpse(df); summary(df)
skimr::skim(df)                     # the best single overview available in R

df %>% summarise(across(everything(), ~ mean(is.na(.)))) %>%
       pivot_longer(everything()) %>% arrange(desc(value))   # % missing
df %>% count(programme, sort = TRUE)
sum(duplicated(df))
\end{lstlisting}

\section{Wrangling with dplyr}

\begin{lstlisting}[language=R]
df %>%
  filter(age > 18, !is.na(mark)) %>%
  mutate(passed  = mark >= 40,
         band    = cut(attendance, c(0, 50, 75, 100),
                       labels = c("low", "medium", "high")),
         mark_z  = (mark - mean(mark)) / sd(mark)) %>%
  select(student_id, programme, mark, passed, band) %>%
  arrange(desc(mark))

df %>%
  group_by(programme) %>%
  summarise(n         = n(),
            mean_mark = mean(mark, na.rm = TRUE),
            median    = median(mark, na.rm = TRUE),   # report this if skewed
            iqr       = IQR(mark, na.rm = TRUE),
            pass_rate = mean(passed), .groups = "drop")

# joins -- check the row count, exactly as in pandas
before <- nrow(df)
df <- left_join(df, vle, by = "student_id")     # left, not inner
stopifnot(nrow(df) == before)

# tidy data (Chapter 3): wide -> long
df %>% pivot_longer(starts_with("wk"), names_to = "week", values_to = "attended")
\end{lstlisting}

\section{Descriptive and Inferential Statistics}

\begin{lstlisting}[language=R]
mean(x); median(x); sd(x); var(x); IQR(x)
quantile(x, c(.25, .5, .75))

# outlier fences -- Chapter 5
q <- quantile(x, c(.25, .75)); iqr <- IQR(x)
x[x < q[1] - 1.5*iqr | x > q[2] + 1.5*iqr]

t.test(mark ~ group, data = df)                  # two-sample
prop.test(c(68, 79), c(100, 100))                # two proportions -- Chapter 9
cor(df$attendance, df$mark, use = "complete.obs")
cor.test(df$attendance, df$mark)

# effect size -- report WITH the p-value, never instead of it
effectsize::cohens_d(mark ~ group, data = df)

# power -- run BEFORE the study, not after (Chapter 10)
pwr::pwr.2p.test(h = pwr::ES.h(0.78, 0.68), power = 0.8, sig.level = 0.05)
\end{lstlisting}

\section{Modelling}

\begin{lstlisting}[language=R]
# --- base R: excellent diagnostics, which is why R is kept for this ----
m <- lm(mark ~ attendance + quizzes + days_absent, data = df)
summary(m)                     # coefficients, R^2, F
par(mfrow = c(2, 2)); plot(m)  # residuals vs fitted, Q-Q, scale-location, leverage
                               # -- the four plots of Chapter 12, free
confint(m)
broom::tidy(m); broom::glance(m)

g <- glm(passed ~ attendance + quizzes, data = df, family = binomial)
exp(coef(g))                   # odds ratios -- Chapter 13

# --- tidymodels: the pipeline equivalent of scikit-learn --------------
split <- initial_split(df, prop = 0.8, strata = passed)   # stratified
tr <- training(split); te <- testing(split)

rec <- recipe(passed ~ ., data = tr) %>%
  step_impute_median(all_numeric_predictors()) %>%
  step_normalize(all_numeric_predictors()) %>%
  step_dummy(all_nominal_predictors())        # fitted on TRAIN only

mod <- logistic_reg() %>% set_engine("glm")
wf  <- workflow() %>% add_recipe(rec) %>% add_model(mod)

folds <- vfold_cv(tr, v = 5, strata = passed)
fit_resamples(wf, folds, metrics = metric_set(pr_auc, roc_auc)) %>%
  collect_metrics()

final <- fit(wf, tr)
augment(final, te) %>% pr_auc(truth = passed, .pred_TRUE)
augment(final, te) %>% conf_mat(truth = passed, estimate = .pred_class)
\end{lstlisting}

\section{Plotting with ggplot2}

\begin{lstlisting}[language=R]
ggplot(df, aes(mark)) + geom_histogram(bins = 30)

ggplot(df, aes(programme, mark)) +
  geom_boxplot() + geom_jitter(width = .15, alpha = .25)   # show the points too

ggplot(df, aes(attendance, mark, colour = programme)) +
  geom_point(alpha = .6) +
  geom_smooth(method = "lm") +
  facet_wrap(~ programme) +          # per group: guards against Simpson's paradox
  labs(title = "Mark against attendance, by programme",
       x = "attendance (%)", y = "final mark") +
  theme_minimal()

GGally::ggpairs(select(df, where(is.numeric)))   # the whole EDA pass in one line

ggsave("fig.png", width = 8, height = 5, dpi = 150)
\end{lstlisting}

\section{Quarto}

\begin{lstlisting}[language=R]
# report.qmd
# ---
# title: "Analysis"
# format: html
# execute:
#   echo: false
#   warning: false
# ---
quarto::quarto_render("report.qmd")
\end{lstlisting}

\begin{conceptbox}[When to Reach for R Rather Than Python]
Use R for exploratory statistics, regression diagnostics (\texttt{plot(m)} gives
all four residual plots of \chref{regression} in one line), classical hypothesis
testing, and any report where the narrative and the analysis belong in one Quarto
document. Use Python for production pipelines, deep learning and anything that must
be deployed as a service. Both are professional choices; using the wrong one for
the task is not.
\end{conceptbox}

% \texorpdfstring keeps the maths out of the PDF bookmark, which is a plain
% string and cannot hold a math shift.
\section{\texorpdfstring{Python $\leftrightarrow$ R Translation}{Python <-> R Translation}}

\rowcolors{2}{white}{lightbg}
\begin{longtable}{@{}L{6.6cm}L{7.4cm}@{}}
\toprule
\rowcolor{primary!15}
\textbf{pandas / scikit-learn} & \textbf{tidyverse / tidymodels} \\
\midrule
\endhead
\texttt{pd.read\_csv("f.csv")} & \texttt{read\_csv("f.csv")} \\
\texttt{df.shape} & \texttt{dim(df)} \\
\texttt{df.head()} & \texttt{head(df)} / \texttt{glimpse(df)} \\
\texttt{df[df.a > 5]} & \texttt{filter(df, a > 5)} \\
\texttt{df[["a","b"]]} & \texttt{select(df, a, b)} \\
\texttt{df.assign(c=df.a+df.b)} & \texttt{mutate(df, c = a + b)} \\
\texttt{df.groupby("g").mean()} & \texttt{df \%>\% group\_by(g) \%>\% summarise(...)} \\
\texttt{df.merge(x, how="left")} & \texttt{left\_join(df, x)} \\
\texttt{df.sort\_values("a")} & \texttt{arrange(df, a)} \\
\texttt{pd.melt / .pivot} & \texttt{pivot\_longer / pivot\_wider} \\
\texttt{train\_test\_split(stratify=y)} & \texttt{initial\_split(strata = y)} \\
\texttt{Pipeline / ColumnTransformer} & \texttt{recipe() \%>\% step\_*()} \\
\texttt{cross\_val\_score} & \texttt{fit\_resamples()} \\
\texttt{classification\_report} & \texttt{conf\_mat() / metric\_set()} \\
\bottomrule
\end{longtable}
